\section{Schema Generală a Aplicației}

Aplicația StyleGenAI urmărește o arhitectură distribuită, separând responsabilitățile de interfață (responsabilitatea dispozitivului mobil) de logica complexă de procesare a datelor (responsabilitatea serverului backend). Structura monorepo permite accesul unitar la întregul cod sursă.

\subsection{Organizarea Directoarelor}
Codul sursă este structurat în două secțiuni majore, reflectând arhitectura client-server:

\begin{itemize}
    \item \textbf{Backend Codebase}: Structurat modular pentru a izola funcționalitățile distincte. Fiecare modul (Auth, Styling, Trends) are propriul director și logica encapsulată.
    \item \textbf{Frontend Codebase}: Utilizează structura de directoare standard Expo, organizată pe ecrane (`screens`) și componente reutilizabile (`components`).
\end{itemize}

Modelul de organizare facilitează atât dezvoltarea individuală a componentelor, cât și integrarea continuă (CI/CD), permițând testarea izolată a funcționalităților înainte de integrarea în fluxul principal al aplicației.

\subsection{Fluxul de Date la Nivel de Modul}
Implementarea urmează fluxul MVC (Model-View-Controller) adaptat pentru API-uri REST:
\begin{enumerate}
    \item \textbf{Controller (Rute API)}: Primește cererea HTTP, validează input-ul și apelează serviciul corespunzător.
    \item \textbf{Service Layer}: Conține logica de business (ex: algoritmul de generare outfit, apelul către modelul AI). Aceasta este inima aplicației.
    \item \textbf{Data Access Layer (Repository)}: Interacționează direct cu baza de date SQL pentru operații CRUD, abstraizând interogările complexe.
\end{enumerate}

\begin{figure}[h]
    \centering
    \includegraphics[width=0.8\textwidth]{diagrama_arhitectura_mvc}
    \caption{Fluxul de date MVC în arhitectura StyleGenAI: De la Controller la Baza de Date via Service Layer.}
    \label{fig:mvc_architecture}
\end{figure}

Această separare asigură un cod testabil și ușor de întreținut, permițând modificarea implementării interne a unui serviciu fără a afecta restul aplicației.
